\section{Conclusion}
\label{sec:conc}

This paper presented an AVX-512 optimized implementation for all six
variants of AIMer, a signature scheme selected as a final algorithm
in the KpqC competition. We parallelized binary-field multiplication across four parties per 512-bit register using VPCLMULQDQ. We also handled conditional XOR accumulation in the linear layer using VPTERNLOGQ. In addition, we accelerated Keccak using AVX-512VL. As a result, our implementation achieved a $1.60$--$1.84\times$ speedup over AVX2 for signing.

We further decomposed this improvement into the contributions of field arithmetic and Keccak. The results show that, in the 128-bit and 192-bit variants, about 86\% of the improvement comes from Keccak. Only in the 256-bit variants does the field contribution increase to about 42\%. These contribution ratios are specific to the selected AVX2 baseline. For the evaluated transition from AVX2 to AVX-512, Keccak accounts for most of the remaining improvement after field arithmetic has already been vectorized.

Our measurements are limited to a single Tiger Lake CPU with Turbo Boost disabled and one pinned core; other AVX-512 microarchitectures, frequency throttling, and multicore behavior remain unexamined.

Future work will evaluate these factors across other AVX-512
microarchitectures and extend the decomposition to ARM SVE and NEON. Our AVX2-to-AVX-512 results suggest that symmetric primitives remain important optimization targets once field arithmetic has already been vectorized.
